% App.I — XDC Pin Map QMTech XC7A100T / XC7A200T
% φ² + φ⁻² = 3 · Trinity S³AI
% Author: Dmitrii Vasilev (ORCID 0009-0008-4294-6159)

\chapter*{Appendix I: XDC Pin Map — QMTech XC7A100T/200T}

\begin{figure}[H]
\centering
\includegraphics[width=0.92\textwidth,keepaspectratio]{app-i-xdc-pin-map.png}
\end{figure}


\addcontentsline{toc}{chapter}{Appendix I: XDC Pin Map}
\label{app:xdc-pin-map}

\textbf{Anchor invariant:} \(\varphi^2 + \varphi^{-2} = 3\).

\section*{I.1 Board Overview}

The QMTech XC7A100T Starter Kit (actual device: XC7A200T, see Appendix~F)
exposes FPGA I/O through a 2×40-pin expansion header (P1/P2) and an on-board
50~MHz oscillator. The following XDC constraints were validated against the
schematic revision~2.1 and confirmed functional via UART loopback test
(Chapter~\ref{ch:32}).

\section*{I.2 Clock}

\begin{table}[h]
\centering
\begin{tabular}{llll}
\toprule
\textbf{Signal} & \textbf{FPGA Pin} & \textbf{Standard} & \textbf{Note} \\
\midrule
clk\_50mhz & U18 & LVCMOS33 & On-board oscillator \\
\bottomrule
\end{tabular}
\caption{Clock constraint.}
\end{table}

\begin{verbatim}
create_clock -period 20.000 -name clk50 [get_ports clk_50mhz]
set_property PACKAGE_PIN U18 [get_ports clk_50mhz]
set_property IOSTANDARD LVCMOS33 [get_ports clk_50mhz]
\end{verbatim}

\section*{I.3 UART (P1 Header)}

\begin{table}[h]
\centering
\begin{tabular}{lllll}
\toprule
\textbf{Signal} & \textbf{Pin} & \textbf{Header} & \textbf{Std} & \textbf{Dir} \\
\midrule
uart\_tx & D20 & P1-11 & LVCMOS33 & Out \\
uart\_rx & E19 & P1-13 & LVCMOS33 & In \\
\bottomrule
\end{tabular}
\caption{UART pins, validated via \filepath{picocom -b 115200 /dev/cu.usbserial-10}.}
\end{table}

\section*{I.4 LED}

\begin{table}[h]
\centering
\begin{tabular}{llll}
\toprule
\textbf{Signal} & \textbf{Pin} & \textbf{Std} & \textbf{Note} \\
\midrule
led[0] & R14 & LVCMOS33 & Active HIGH \\
led[1] & P14 & LVCMOS33 & Active HIGH \\
led[2] & N16 & LVCMOS33 & Active HIGH \\
led[3] & M16 & LVCMOS33 & Active HIGH \\
\bottomrule
\end{tabular}
\caption{On-board LEDs.}
\end{table}

\section*{I.5 JTAG WiFi Bridge (ESP32 → FPGA)}

The JTAG signals are brought out via P2 header to the ESP32 XVC bridge.
Pinout confirmed functional (IDCODE~\texttt{0x13631093}, STAT~\texttt{0x401079FC}).

\begin{table}[h]
\centering
\begin{tabular}{lllll}
\toprule
\textbf{JTAG} & \textbf{ESP32 GPIO} & \textbf{P2 Pin} & \textbf{FPGA Pin} & \textbf{Wire} \\
\midrule
TMS & IO18 & P2-3 & — & Yellow \\
TCK & IO19 & P2-5 & — & Blue \\
TDI & IO23 & P2-7 & — & Green \\
TDO & IO35 & P2-9 & — & White \\
GND & GND & P2-2 & — & Black \\
\bottomrule
\end{tabular}
\caption{JTAG wiring. IO35 on ESP32 is input-only — correct for TDO (FPGA output).}
\end{table}

\section*{I.6 Full XDC Listing}

\begin{verbatim}
## Clock
set_property PACKAGE_PIN U18 [get_ports clk_50mhz]
set_property IOSTANDARD LVCMOS33 [get_ports clk_50mhz]
create_clock -period 20.000 -name clk50 [get_ports clk_50mhz]

## UART
set_property PACKAGE_PIN D20 [get_ports uart_tx]
set_property IOSTANDARD LVCMOS33 [get_ports uart_tx]
set_property PACKAGE_PIN E19 [get_ports uart_rx]
set_property IOSTANDARD LVCMOS33 [get_ports uart_rx]

## LEDs
set_property PACKAGE_PIN R14 [get_ports {led[0]}]
set_property PACKAGE_PIN P14 [get_ports {led[1]}]
set_property PACKAGE_PIN N16 [get_ports {led[2]}]
set_property PACKAGE_PIN M16 [get_ports {led[3]}]
set_property IOSTANDARD LVCMOS33 [get_ports led]

## Timing false paths (async outputs)
set_false_path -to [get_ports uart_tx]
set_false_path -to [get_ports led]
\end{verbatim}
